% source: https://tex.stackexchange.com/a/763093 (question 763091, score 6, block 1)
\documentclass[border=5pt]{standalone}% compile with lualatex only
\usepackage[svgnames]{xcolor}
\usepackage[3d]{luadraw}%https://github.com/pfradin/luadraw
\usepackage{fourier-otf}
%https://tex.stackexchange.com/questions/763091/how-to-make-3d-axes-pass-above-below-a-surface-in-pgfplots
\begin{document}
% first version
\begin{luadraw}{name=axes1}
local ld = luadraw
local pt3d = ld.pt3d
local M = pt3d.M
local g = ld.graph3d:new{ window3d={-5,5,-5,6,-4,4}, window={-6,6,-6,5}, 
    viewdir={-40,70}, size={12,12}}
local f = function(x,y) return 10*y/((x^2+1)*(y^2+2*y+4)) end
local S, wall = ld.cartesian3d(f, -4,4,-4,4,{35,25},"xy")
g:Dscene3d(
    g:addWall(wall), -- partitioning with planes x= and y=
    g:addFacet(S, {usepalette={ld.palRainbow,"z"}, edge=true, edgecolor="gray", 
        edgewidth=1,opacity=0.8}),
    g:addAxes(M(0,0,0.05), {arrows=1})
)
g:Show()
\end{luadraw}

%second version
\begin{luadraw}{name=axes2}
local ld = luadraw
local pt3d = ld.pt3d
local M = pt3d.M
local g = ld.graph3d:new{ window3d={-5,5,-5,5,-4,4}, window={-8.5,8,-7,7}, 
    viewdir={-40,70}, size={12,12}}
local f = function(x,y) return 10*y/((x^2+1)*(y^2+2*y+4)) end
local S = ld.cartesian3d(f, -4,4,-4,4,{35,25})
g:Dboxaxes3d( {grid=true, gridcolor="gray", fillcolor="lightgray"} )
g:Dfacet( S, {usepalette={ld.palRainbow,"z"}, edgecolor="gray", 
        edgewidth=1,opacity=0.6})
g:Show()
\end{luadraw}
\end{document}
